% ==============================================================================
% IEEE Transactions on Wireless Communications (TWC) — LaTeX Template
% Target: IEEE TWC
% Title: (Provisional)
% ==============================================================================
\documentclass[journal]{IEEEtran}

% ---- Packages ----------------------------------------------------------------
\usepackage{cite}
\usepackage{amsmath,amssymb,amsfonts}
\usepackage{algorithmic}
\usepackage{algorithm}
\usepackage{graphicx}
\usepackage{textcomp}
\usepackage{xcolor}
\usepackage{multirow}
\usepackage{booktabs}
\usepackage{url}
\usepackage{hyperref}
\usepackage{subcaption}
\usepackage{bm}

% ---- Math Operators ----------------------------------------------------------
\DeclareMathOperator*{\argmin}{arg\,min}
\DeclareMathOperator*{\argmax}{arg\,max}
\newcommand{\norm}[1]{\left\| #1 \right\|}

% ==============================================================================
\begin{document}

\title{AoI-Aware V2I Uplink Scheduling via Deep Reinforcement Learning with Hybrid Action Space}

\author{Youngju~Nam,~Dick~Mugerwa,~Yongje~Shin,~Jeong-Hun~Kim,~and~Euisin~Lee,~\IEEEmembership{}%
\thanks{This work was supported by the National Research Foundation of Korea (NRF) grant funded by the Korean government Ministry of Science and ICT (MSIT) under Grant No. 2022R1A2C1010121.
\textit{(Corresponding author: Euisin Lee)}}%
\thanks{Youngju Nam is with the School of Software, Kunsan National University, Gunsan 54150, Republic of Korea (e-mail: imnyj@kunsan.ac.kr).}%
\thanks{Dick Mugerwa is with Kyambogo University, Kampala, Uganda (e-mail: dmugerwa@kyu.ac.ug).}%
\thanks{Yongje Shin is with the School of Smart IT, Semyung University, Jecheon 27136, Republic of Korea (e-mail: dudujr@semyung.ac.kr).}%
\thanks{Jeong-Hun Kim is with the Department of Computer Science and Information Engineering, Kunsan National University, Gunsan 54150, Republic of Korea (e-mail: kimjh@kunsan.ac.kr).}%
\thanks{Euisin Lee is with the School of Information and Communication Engineering, Chungbuk National University, Cheongju 28644, Republic of Korea (e-mail: eslee@chungbuk.ac.kr).}%
}

\markboth{IEEE Transactions on Wireless Communications}%
{Nam \MakeLowercase{\textit{et al.}}: AoI-Aware V2I Uplink Scheduling via DRL}

\maketitle

% ==============================================================================
% ABSTRACT
% ==============================================================================
\begin{abstract}
In vehicle-to-infrastructure (V2I) networks, roadside units (RSUs) maintain real-time awareness of surrounding vehicles by receiving periodic mobility updates.
However, indiscriminate update scheduling wastes uplink resources and exacerbates channel congestion, while infrequent updates degrade the freshness of information.
This paper formulates the V2I uplink scheduling problem as a semi-Markov decision process (SMDP) that jointly optimizes the update interval, transmit power, and subchannel selection for each vehicle.
We propose a deep reinforcement learning (DRL) framework with a hybrid action space that combines continuous control of the update interval $\Delta$ and transmit power $p$ with discrete subchannel selection.
The reward function penalizes estimation error, power consumption, channel congestion, and redundant updates, where the weighting coefficients are jointly optimized via Optuna hyperparameter search.
We build a genuine simulation environment using SUMO microscopic traffic simulator integrated with a 5.9\,GHz Rayleigh fading channel model, and evaluate the proposed approach against nine baseline methods including three foundational RL algorithms and six recent AoI-aware scheduling schemes.
Simulation results demonstrate that the proposed method achieves a favorable trade-off between information freshness and resource efficiency under varying traffic densities.
\end{abstract}

\begin{IEEEkeywords}
Age of information, vehicle-to-infrastructure, uplink scheduling, deep reinforcement learning, hybrid action space, semi-Markov decision process.
\end{IEEEkeywords}

% ==============================================================================
% I. INTRODUCTION
% ==============================================================================
\section{Introduction}
\IEEEPARstart{T}{he} proliferation of connected vehicles and intelligent transportation systems (ITS) has created an increasing demand for timely and accurate mobility information at roadside infrastructure.
In vehicle-to-infrastructure (V2I) networks, roadside units (RSUs) maintain a local awareness table that tracks the position, velocity, and heading of surrounding vehicles.
The freshness of this information, commonly quantified by the Age of Information (AoI)~\cite{kaul2012real}, is critical for safety-critical applications such as cooperative collision avoidance, intersection management, and autonomous driving support.

% --- Problem Statement ---
Managing V2I uplink updates through cellular networks incurs significant cost.
Instead, RSUs can selectively schedule updates only for vehicles whose mobility information is deemed stale or unreliable.
This selective scheduling introduces a fundamental trade-off: frequent updates improve information freshness but consume more power and increase channel congestion, while infrequent updates conserve resources but allow the RSU's estimate to diverge from reality.

% --- Limitations of Prior Work ---
Existing AoI-aware scheduling approaches~\cite{mlika2022deep, qi2025deep, zhang2025drl} typically treat the update decision as a simple binary choice or optimize only a single dimension (e.g., update interval or power) in isolation.
Moreover, most prior works adopt simplified channel models that do not capture the stochastic nature of V2I uplink transmissions under Rayleigh fading and multi-user interference.

% --- Our Approach ---
In this paper, we formulate the V2I uplink scheduling problem as a semi-Markov decision process (SMDP) and propose a deep reinforcement learning framework with a \emph{hybrid action space}.
The agent jointly decides three control variables for each vehicle: (i) a continuous update interval $\Delta \in [\Delta_{\min}, \Delta_{\max}]$, (ii) a continuous transmit power $p \in [p_{\min}, p_{\max}]$, and (iii) a discrete subchannel index $\mathrm{ch} \in \{0, 1, \ldots, C{-}1\}$.
The reward function incorporates four penalty terms---estimation error, power consumption, channel congestion, and redundant updates---each normalized to $[0, 1]$ and weighted by coefficients $w_1 \sim w_4$ that are optimized via Optuna hyperparameter search.

% --- Contributions ---
The main contributions of this paper are summarized as follows:
\begin{itemize}
    \item We formulate the AoI-aware V2I uplink scheduling problem as an SMDP with a hybrid action space that jointly optimizes the update interval, transmit power, and subchannel selection.
    \item We design a four-term reward function with Min-Max normalization and Optuna-optimized weighting coefficients that captures the multi-dimensional trade-off between information freshness, resource consumption, and channel congestion.
    \item We implement a genuine simulation environment based on the SUMO microscopic traffic simulator with a physically accurate 5.9\,GHz Rayleigh fading channel model, and conduct rigorous evaluation against nine baseline methods.
    \item We introduce an Act/Rest dual-model hot-swap training architecture that enables continuous inference during background model updates, ensuring hardware feasibility for real-time deployment.
\end{itemize}

The remainder of this paper is organized as follows. Section~\ref{sec:related} reviews related work. Section~\ref{sec:system} presents the system model. Section~\ref{sec:problem} formulates the SMDP problem. Section~\ref{sec:proposed} describes the proposed DRL framework. Section~\ref{sec:evaluation} presents the simulation setup and results. Section~\ref{sec:conclusion} concludes the paper.

% ==============================================================================
% II. RELATED WORK
% ==============================================================================
\section{Related Work}
\label{sec:related}

\subsection{Age of Information in Vehicular Networks}
The concept of AoI was originally introduced for status update systems~\cite{kaul2012real} and has since been extensively studied in vehicular networks.

Mlika and Cherkaoui~\cite{mlika2022deep} applied DDPG to minimize AoI in cellular V2X communications, demonstrating the feasibility of DRL-based scheduling for vehicular scenarios.
Qi \emph{et~al.}~\cite{qi2025deep} proposed a SAC-based resource allocation scheme for RIS-aided IoV networks, incorporating AoI awareness into the reward design.
Arani \emph{et~al.}~\cite{arani2026resilient} developed a resilient AoI-aware optimization framework using DRL for intelligent transportation systems.
Zhang \emph{et~al.}~\cite{zhang2025drl} jointly optimized AoI and energy consumption in C-V2X enabled IoV using DRL.

\subsection{Resource Allocation in V2X Networks}
Azizi \emph{et~al.}~\cite{azizi2024efficient} proposed a multi-agent RL mechanism for AoI-aware resource management in VLC-V2X networks.
Lin \emph{et~al.}~\cite{lin2025optimization} optimized spectrum resource allocation for vehicle platoons using DRL.

\subsection{Hybrid Action Spaces in DRL}
While prior works predominantly adopt either fully continuous or fully discrete action spaces, our work jointly handles continuous (interval, power) and discrete (subchannel) variables through a hybrid action space architecture.

% ==============================================================================
% III. SYSTEM MODEL
% ==============================================================================
\section{System Model}
\label{sec:system}

\subsection{Network Architecture}
We consider a V2I network where RSUs are deployed at intersections of an urban grid road network.
Each RSU has a communication range of $R_{\mathrm{RSU}} = 300$\,m operating at a carrier frequency of 5.9\,GHz, consistent with the C-V2X sidelink specification.
The uplink bandwidth of 20\,MHz is divided into $C = 4$ subchannels.
Vehicles enter and exit the RSU's coverage area as they traverse the road network, which is simulated using the SUMO microscopic traffic simulator with a step resolution of 0.1\,s.

\subsection{RSU Awareness Table}
Each RSU maintains an awareness table that stores, for each registered vehicle $i$, the last received mobility state $(\mathbf{x}_i, \mathbf{v}_i, \tau_i)$, where $\mathbf{x}_i$ is the position, $\mathbf{v}_i$ is the velocity, and $\tau_i$ is the timestamp of the last successful update.
Between updates, the RSU extrapolates the vehicle's position using a constant-velocity model:
\begin{equation}
    \hat{\mathbf{x}}_i(t) = \mathbf{x}_i + \mathbf{v}_i \cdot (t - \tau_i).
    \label{eq:extrapolation}
\end{equation}

The estimation error at time $t$ is defined as:
\begin{equation}
    e_i(t) = \| \hat{\mathbf{x}}_i(t) - \mathbf{x}_i^{\mathrm{true}}(t) \|_2,
    \label{eq:error}
\end{equation}
where $\mathbf{x}_i^{\mathrm{true}}(t)$ is the true position obtained from the SUMO simulator.

\subsection{Wireless Channel Model}
The uplink transmission from vehicle $i$ to the RSU follows a path loss model with Rayleigh fading.
The received power at distance $d_i$ is:
\begin{equation}
    P_{\mathrm{rx},i} = P_{\mathrm{tx},i} - \mathrm{PL}(d_i),
\end{equation}
where the path loss is given by:
\begin{equation}
    \mathrm{PL}(d) = 20\log_{10}\!\left(\frac{4\pi f_c}{c}\right) + 10 \alpha \log_{10}(d) \quad [\mathrm{dB}],
\end{equation}
with carrier frequency $f_c = 5.9$\,GHz and path loss exponent $\alpha = 2.3$.

For vehicles transmitting on the same subchannel, the success probability under Rayleigh fading is given by the closed-form expression:
\begin{equation}
    P_{\mathrm{succ},i} = \exp\!\left(-\frac{\gamma_{\mathrm{th}} N_0}{S_i}\right) \prod_{k \neq i} \frac{1}{1 + \gamma_{\mathrm{th}} I_k / S_i},
    \label{eq:sinr}
\end{equation}
where $\gamma_{\mathrm{th}}$ is the SINR threshold, $N_0$ is the noise power, $S_i$ is the received signal power from vehicle $i$, and $I_k$ is the interference from co-channel vehicle $k$.

\subsection{Vehicle Events}
A vehicle's lifecycle within an RSU cell consists of three events:
\begin{itemize}
    \item \textbf{E1 (Entry)}: Vehicle enters the RSU's coverage, performs initial registration, and receives its first update schedule.
    \item \textbf{E2 (Scheduled Update)}: At the scheduled time, the vehicle transmits its current state. The transmission succeeds or fails according to~\eqref{eq:sinr}.
    \item \textbf{E3 (Exit)}: Vehicle leaves the coverage area and is deregistered.
\end{itemize}

% ==============================================================================
% IV. PROBLEM FORMULATION
% ==============================================================================
\section{Problem Formulation}
\label{sec:problem}

\subsection{Semi-Markov Decision Process}
We model the scheduling problem as an SMDP, where the decision epochs are not uniformly spaced but occur at each vehicle's scheduled update time.

\subsubsection{State Space}
The state observation for vehicle $i$ at decision epoch $t$ is an 18-dimensional vector:
\begin{equation}
    \mathbf{s}_i(t) = \left[ v_{\mathrm{pos}}, d_{\mathrm{rsu}}, v_{\mathrm{heading}}, v_{\mathrm{vel}}, \mathrm{tls}_{\mathrm{state}}, \mathrm{tls}_{\mathrm{dist}}, n_{\mathrm{queue}}, n_{\mathrm{active}}, \ldots \right]^{\!\top}\!,
\end{equation}
which includes:
\begin{itemize}
    \item Vehicle kinematics: position $v_{\mathrm{pos}}$, distance to RSU $d_{\mathrm{rsu}}$, heading $v_{\mathrm{heading}}$ (approaching/leaving), velocity $v_{\mathrm{vel}}$.
    \item Traffic signal context: signal state $\mathrm{tls}_{\mathrm{state}}$ (R/Y/G), distance to stop line $\mathrm{tls}_{\mathrm{dist}}$.
    \item Congestion indicators: number of queued vehicles ahead $n_{\mathrm{queue}}$, number of active vehicles in range $n_{\mathrm{active}}$.
    \item Historical context: information from previously updated neighboring vehicles stored in the RSU table.
\end{itemize}

\subsubsection{Action Space}
The hybrid action for vehicle $i$ is a tuple:
\begin{equation}
    \mathbf{a}_i = (\Delta_i,\; p_i,\; \mathrm{ch}_i),
\end{equation}
where:
\begin{itemize}
    \item $\Delta_i \in [\Delta_{\min}, \Delta_{\max}]$: next update interval (continuous). The upper bound $\Delta_{\max}$ is set to the maximum red-phase duration of traffic lights in the simulation.
    \item $p_i \in [p_{\min}, p_{\max}] = [10, 23]$\,dBm: transmit power (continuous), bounded by the 3GPP power class~3 UE maximum.
    \item $\mathrm{ch}_i \in \{0, 1, \ldots, C{-}1\}$: subchannel index (discrete).
\end{itemize}

\subsubsection{Reward Function}
The reward is a weighted sum of four normalized penalty terms:
\begin{equation}
    R_t = -\!\left( w_1 \cdot \overline{e_t^2} + w_2 \cdot \overline{P_{\mathrm{tx}}} + w_3 \cdot \overline{C_{\mathrm{freq}}} + w_4 \cdot \mathbb{I}_{\mathrm{redundant}} \right),
    \label{eq:reward}
\end{equation}
where the overline denotes Min-Max normalization to $[0,1]$:
\begin{itemize}
    \item $\overline{e_t^2}$: normalized squared estimation error between the RSU's extrapolated and true vehicle positions.
    \item $\overline{P_{\mathrm{tx}}} = (p - p_{\min}) / (p_{\max} - p_{\min})$: normalized transmit power penalty.
    \item $\overline{C_{\mathrm{freq}}}$: normalized channel congestion indicator based on the channel busy ratio (CBR) and SINR collision rate.
    \item $\mathbb{I}_{\mathrm{redundant}} \in \{0, 1\}$: binary indicator that imposes a penalty when the vehicle's physical state has not changed since the last update (e.g., a stopped vehicle).
\end{itemize}

The weighting coefficients $w_1, w_2, w_3, w_4$ are not heuristically fixed but are included in the Optuna hyperparameter search space, allowing the agent to discover the optimal balance across penalty terms.

\subsection{SMDP Discounting}
Since the time between successive actions for a given vehicle varies, we employ the SMDP discount factor $\gamma^{\Delta_i}$, which accounts for the variable sojourn time in the replay buffer.

% ==============================================================================
% V. PROPOSED METHOD
% ==============================================================================
\section{Proposed DRL Framework}
\label{sec:proposed}

\subsection{Hybrid Action Architecture}
The proposed DRL agent outputs the hybrid action $\mathbf{a}_i = (\Delta_i, p_i, \mathrm{ch}_i)$ through a multi-head neural network.
Two continuous heads produce $\Delta_i$ and $p_i$ via tanh activation with subsequent scaling, while a discrete head selects $\mathrm{ch}_i$ via categorical sampling.

\subsection{Act/Rest Dual-Model Hot-Swap}
To ensure real-time inference feasibility, we adopt an Act/Rest dual-model architecture:
\begin{itemize}
    \item \textbf{Act mode}: The active model serves real-time scheduling decisions while collecting transitions into a replay buffer.
    \item \textbf{Rest mode}: A background model trains on the accumulated transitions using a separate hardware resource.
    \item \textbf{Hot-swap}: When the Rest model's training is complete, it seamlessly replaces the Act model with zero downtime, and the cycle repeats.
\end{itemize}
This architecture ensures that the training workload does not interfere with inference latency, which is critical for time-sensitive V2I scheduling.

\subsection{Hyperparameter Optimization}
We employ Optuna for hyperparameter optimization (HPO), searching over the reward weights $w_1 \sim w_4$ and model-specific hyperparameters (learning rate, discount factor, etc.).
Each trial trains the agent and evaluates it across multiple seeds and traffic densities.

% ==============================================================================
% VI. PERFORMANCE EVALUATION
% ==============================================================================
\section{Performance Evaluation}
\label{sec:evaluation}

\subsection{Simulation Setup}

\subsubsection{Traffic Simulation}
We use SUMO (Simulation of Urban Mobility) v1.27 to generate realistic vehicular traffic on an urban grid network.
Table~\ref{tab:sim_params} summarizes the key simulation parameters.

\begin{table}[!t]
\centering
\caption{Simulation Parameters}
\label{tab:sim_params}
\begin{tabular}{lc}
\toprule
\textbf{Parameter} & \textbf{Value} \\
\midrule
Carrier frequency & 5.9\,GHz \\
Bandwidth & 20\,MHz \\
Number of subchannels & 4 \\
RSU communication range & 300\,m \\
Transmit power range & [10, 23]\,dBm \\
Path loss exponent & 2.3 \\
SINR threshold & 0\,dB \\
SUMO step length & 0.1\,s \\
Episode length & 2000 steps (200\,s) \\
Number of episodes & 100 (200k steps total) \\
Vehicle density & 15--55\,veh/km \\
\bottomrule
\end{tabular}
\end{table}

\subsubsection{Baseline Methods}
We compare the proposed method against nine baselines:
\begin{itemize}
    \item \textbf{Foundation RL models}: PPO~\cite{ppo}, SAC~\cite{sac}, TD3~\cite{td3}.
    \item \textbf{AoI-aware DRL schemes}: SAC-RIS~\cite{qi2025deep}, DDPG-CV2X~\cite{mlika2022deep}, DDPG-Resilient~\cite{arani2026resilient}, MARL-VLC~\cite{azizi2024efficient}, Platoon-DRL~\cite{lin2025optimization}, DRL-IoV~\cite{zhang2025drl}.
\end{itemize}

All baselines are adapted to the same hybrid action space and trained on the identical SUMO environment for fair comparison.

\subsection{Training Configuration}
Each model is trained for 200,000 environment steps (100 episodes $\times$ 2,000 steps).
Optuna HPO is used to optimize the reward weights and model hyperparameters over multiple trials.
Training progress is monitored via TensorBoard.

\subsection{Evaluation Metrics}
We evaluate the following metrics across varying vehicle densities (15--55\,veh/km) and multiple random seeds:
\begin{itemize}
    \item Mean estimation error ($\overline{e^2}$).
    \item Uplink success rate.
    \item Mean AoI across all vehicles.
    \item Total energy consumption.
    \item Channel utilization efficiency.
\end{itemize}

\subsection{Results and Discussion}

% Placeholder for results
\textit{(Results to be added after the 200,000-step training runs are completed and analyzed.)}

% ==============================================================================
% VII. CONCLUSION
% ==============================================================================
\section{Conclusion}
\label{sec:conclusion}

In this paper, we proposed a DRL-based V2I uplink scheduling framework that jointly optimizes the update interval, transmit power, and subchannel selection through a hybrid action space.
The four-term reward function with Optuna-optimized weights captures the complex trade-off between information freshness, power efficiency, channel congestion, and update redundancy.
The Act/Rest dual-model hot-swap architecture ensures real-time inference feasibility.
Extensive simulations on a genuine SUMO traffic environment with a 5.9\,GHz Rayleigh fading channel model validate the effectiveness of the proposed approach.

% ==============================================================================
% REFERENCES
% ==============================================================================
\bibliographystyle{IEEEtran}
% \bibliography{references}

\begin{thebibliography}{10}

\bibitem{kaul2012real}
S.~Kaul, R.~Yates, and M.~Gruteser, ``Real-time status: How often should one update?'' in \emph{Proc. IEEE INFOCOM}, 2012, pp.~2187--2195.

\bibitem{ppo}
J.~Schulman, F.~Wolski, P.~Dhariwal, A.~Radford, and O.~Klimov, ``Proximal policy optimization algorithms,'' \emph{arXiv preprint arXiv:1707.06347}, 2017.

\bibitem{sac}
T.~Haarnoja, A.~Zhou, P.~Abbeel, and S.~Levine, ``Soft actor-critic: Off-policy maximum entropy deep reinforcement learning with a stochastic actor,'' in \emph{Proc. Int. Conf. Mach. Learn. (ICML)}, 2018, pp.~1856--1865.

\bibitem{td3}
S.~Fujimoto, H.~van~Hoof, and D.~Meger, ``Addressing function approximation error in actor-critic methods,'' in \emph{Proc. Int. Conf. Mach. Learn. (ICML)}, 2018.

\bibitem{qi2025deep}
K.~Qi, Q.~Wu, P.~Fan, N.~Cheng, W.~Chen, J.~Wang, and K.~B.~Letaief, ``Deep-reinforcement-learning-based AoI-aware resource allocation for RIS-aided IoV networks,'' \emph{IEEE Trans. Veh. Technol.}, vol.~74, no.~1, 2025, doi:~10.1109/TVT.2024.3452790.

\bibitem{mlika2022deep}
Z.~Mlika and S.~Cherkaoui, ``Deep deterministic policy gradient to minimize the age of information in cellular V2X communications,'' \emph{IEEE Trans. Intell. Transp. Syst.}, vol.~23, no.~12, pp.~23597--23612, 2022, doi:~10.1109/TITS.2022.3190799.

\bibitem{arani2026resilient}
A.~H.~Arani, H.~Saeedi, S.~Norouzi, A.~Nouruzi, N.~Zorba, and H.~Yanikomeroglu, ``A resilient AoI-aware optimization framework for intelligent transportation systems using deep reinforcement learning,'' \emph{IEEE Open J. Commun. Soc.}, 2026, doi:~10.1109/OJCOMS.2026.3707734.

\bibitem{azizi2024efficient}
M.~Azizi, F.~Zeinali, M.~R.~Mili, and S.~Shokrollahi, ``Efficient AoI-aware resource management in VLC-V2X networks via multi-agent RL mechanism,'' \emph{IEEE Trans. Veh. Technol.}, vol.~73, no.~9, pp.~14009--14014, 2024, doi:~10.1109/TVT.2024.3392738.

\bibitem{lin2025optimization}
Z.~Lin, H.~Pan, and Y.~Wang, ``Optimization of spectrum resource allocation for vehicle platoon in V2X networks based on deep reinforcement learning,'' \emph{IEEE Trans. Veh. Technol.}, vol.~75, no.~6, pp.~11512--11527, 2025, doi:~10.1109/TVT.2025.3643923.

\bibitem{zhang2025drl}
Z.~Zhang, Q.~Wu, P.~Fan, N.~Cheng, W.~Chen, and K.~B.~Letaief, ``DRL-based optimization for AoI and energy consumption in C-V2X enabled IoV,'' \emph{IEEE Trans. Green Commun. Netw.}, vol.~9, no.~4, pp.~2144--2159, 2025, doi:~10.1109/TGCN.2025.3531902.

\end{thebibliography}

\end{document}
